\Subhead{2.7.1}{Метод Кирхгофа. Функции Грина. Условия излучения}

Монохроматическое скалярное поле в однородной среде удовлетворяет $(\Delta+k^2)U=0$. Исходящая функция Грина
\begin{equation}
 G(\bm r,\bm r')=\frac{\e^{\ii k|\bm r-\bm r'|}}{4\pi|\bm r-\bm r'|},
 \quad (\Delta+k^2)G=-\delta(\bm r-\bm r')
\end{equation}
выбирается условием Зоммерфельда $r(\partial_rU-\ii kU)\to0$ при $r\to\infty$. Оно исключает волну, приходящую из бесконечности. Вторая формула Грина даёт для замкнутой поверхности с внешней нормалью
\begin{equation}
 U(\bm r)=\oint_{\partial V}\left[G\partial_nU-U\partial_nG\right]\dd S'
\end{equation}
при наблюдении внутри области без источников. Знак зависит от ориентации нормали и принятого знака дельта-источника; здесь они заданы совместно.

Приближение Кирхгофа полагает поле и его производную в отверстии равными падающей волне, а на непрозрачном экране — нулевыми. Эти граничные предположения у края не являются точным решением; метод хорошо работает для большого относительно $\lambda$ отверстия вдали от тонкой краевой области. Векторная электромагнитная задача требует также поляризации и проводящих граничных условий. Физический смысл интеграла — интерференция вторичных волн, а не независимое излучение каждой точки с произвольной фазой.
